\begin{abstract}
Multimodal attention systems assume that fusing more data streams improves recognition. We test this with gTAO-Net, one model in which each design choice is a switch: gated versus mean-pooled fusion, a linear input-to-head link, modality dropout, and three ordinal objectives (CORAL, a CORE-style hybrid and a soft multi-rater target). All streams share a training-free Tchebichef-moment front-end. On a DIPSER slice of 20 students, none of the 24 fusion configurations exceeds the best single stream (face mesh with the soft target, QWK $0.256\pm0.068$, against $0.124\pm0.071$ for the best fusion), and the gate is not significantly different from mean pooling ($p=0.68$). Only the input-to-head link ($p=0.002$) and the soft target ($p=0.023$) help, and fusion yields no meaningful gain on DAiSEE or EngageNet. On the full DIPSER release (\fcPersons{} students, \fcWindows{} windows) the fusion result replicates (\fcFusionsAbove{} fusions above it), but the sensor signal is weak (best QWK \fcBestSingleQWK) and largely shared by the class. Hyper-parameter search, richer and local moments, a face crop, frozen pretrained face encoders, and temporal and class context do not improve on it. A label probe every 32\,s raises accuracy to \fcProbeHalfAcc\% on an unseen class, but repeating the probe already gives \fcPersistHalfAcc\%. The evaluation protocol has a larger effect than any of these choices: the same model obtains QWK \fcLeakRandQWK{} on random windows, \fcLeakPersonQWK{} on held-out students and \fcLogoOwnQWK{} on a held-out class, while its accuracy remains close to the majority rate (\fcMajority\%).
\end{abstract}

\begin{keywords}
Data leakage, DIPSER, evaluation protocol, multimodal fusion, multi-rater labels, ordinal regression, student attention, Tchebichef moments.
\end{keywords}
